\documentclass[11pt,a4paper]{article}

% --- Packages ---
\usepackage[a4paper, margin=2.5cm]{geometry}
\usepackage{amsmath, amssymb, amsthm}
\usepackage{graphicx}
\usepackage{xcolor}
\usepackage{booktabs}
\usepackage{enumitem}
\usepackage{microtype}
\usepackage{titlesec}
\usepackage{subcaption}
\usepackage[skins, breakable]{tcolorbox}
\usepackage{tikz}
\usetikzlibrary{arrows.meta, positioning, shapes.geometric, fit, backgrounds}
\usepackage[colorlinks=true, linkcolor=blue!60!black,
            citecolor=blue!60!black, urlcolor=blue!70!black]{hyperref}
\usepackage[ruled,vlined,linesnumbered]{algorithm2e}
\usepackage{pgfplots}
\pgfplotsset{compat=1.18}

% --- Color palette ---
\definecolor{accentblue}{RGB}{25, 80, 160}
\definecolor{accentgray}{RGB}{90, 95, 105}
\definecolor{boxbg}{RGB}{240, 245, 255}
\definecolor{boxborder}{RGB}{25, 80, 160}
\definecolor{warnbg}{RGB}{255, 248, 235}
\definecolor{warnborder}{RGB}{200, 130, 30}

% --- Section formatting ---
\titleformat{\section}{\large\bfseries\color{accentblue}}{{\color{accentblue}\thesection.}}{0.5em}{}[\color{accentblue}\titlerule]
\titleformat{\subsection}{\normalsize\bfseries\color{accentgray}}{\thesubsection.}{0.5em}{}
\titleformat{\subsubsection}{\normalsize\itshape\color{accentgray}}{\thesubsubsection.}{0.5em}{}

% --- Tcolorbox styles ---
\newtcolorbox{keybox}[1]{
  enhanced, breakable,
  colback=boxbg, colframe=boxborder,
  fonttitle=\bfseries\color{white},
  coltitle=white, colbacktitle=boxborder,
  title=#1, arc=4pt, boxrule=0.8pt,
  left=6pt, right=6pt, top=4pt, bottom=4pt
}
\newtcolorbox{notebox}{
  enhanced, colback=warnbg, colframe=warnborder,
  arc=4pt, boxrule=0.8pt, fontupper=\small,
  left=6pt, right=6pt, top=4pt, bottom=4pt
}
\newtcolorbox{riskbox}[1]{
  enhanced, breakable,
  colback=red!4!white, colframe=red!45!black,
  fonttitle=\bfseries\color{white},
  coltitle=white, colbacktitle=red!45!black,
  title=#1, arc=4pt, boxrule=0.8pt, fontupper=\small,
  left=6pt, right=6pt, top=4pt, bottom=4pt
}

% --- Math operators ---
\DeclareMathOperator{\SynOps}{SynOps}
\newcommand{\vth}{V_{\mathrm{th}}}
\newcommand{\vmem}{V_{\mathrm{mem}}}

% =========================================================
\begin{document}

% --- Title block ---
\begin{center}
  {\LARGE\bfseries\color{accentblue}
    Memristor-Inspired Neuromorphic Control for Robotics:\\[4pt]
    A Comparative Study of SNN and DNN Architectures\\[4pt]
    with FPGA-Accelerated Synaptic Emulation}\\[1.2em]
  {\large\color{accentgray} Undergraduate Thesis Proposal --- Electrical \& Computer Engineering}\\[0.4em]
\end{center}

\vspace{0.5em}
\hrule height 1.5pt
\vspace{1.2em}

% --- Abstract ---
\begin{keybox}{Abstract}
This thesis proposes a hardware-software co-design framework for neuromorphic robotic
control, combining spiking neural networks (SNNs) with memristor-inspired synaptic
plasticity and FPGA-accelerated weight update logic. A leaky integrate-and-fire (LIF)
SNN with \emph{reward-modulated} spike-timing-dependent plasticity (R-STDP) is deployed
in a simulated obstacle-avoidance task (PyBullet) and evaluated against a trained deep
neural network (DNN) baseline and a reinforcement-learning (RL) policy; a purely Hebbian
STDP variant is retained as an ablation. The FPGA implements the
parallel synaptic conductance update array in hardware, a topology and update rule
that is structurally analogous to the crossbar-array dynamics of physical memristive
devices (a memristor-\emph{inspired} accelerator, not a device-level emulator). The
central hypothesis is that online, reward-gated synaptic adaptation---impossible in a
frozen DNN---enables measurably faster performance recovery when the environment changes
mid-episode, while the FPGA acceleration provides latency advantages on the
weight-update kernel that justify the hardware-software co-design perspective.
Metrics include response latency, synaptic operation count (SynOps), adaptation
recovery time, and robustness under sensor noise.
\end{keybox}

\vspace{1em}

% =========================================================
\section{Motivation and Research Gap}

Modern robotic controllers are overwhelmingly based on deep neural networks trained
offline: once deployed, their weights are frozen. This works well in static
environments but introduces a fundamental brittleness---any distributional shift (new
obstacles, sensor degradation, lighting changes) requires costly retraining. Biological
neural systems solve this through \emph{online synaptic plasticity}: the synapse itself
carries memory and continues adapting throughout the lifetime of the organism.

Memristive devices are the closest physical analog to biological synapses. Their
resistance state encodes a weight, updates depend on the history of current flow (state
retention), the relationship is nonlinear and bounded (saturation), and device
fabrication introduces stochasticity---all hallmarks of biological synaptic behavior
\cite{strukov2008}. An FPGA crossbar-style update array is the most accessible way to
emulate the \emph{structure} of this parallel, in-memory computation without
fabricating physical memristor hardware; it reproduces the topology and update locality
of a crossbar, not the continuous-time device physics.

\begin{keybox}{Core Research Questions}
\begin{enumerate}[label=\textbf{RQ\arabic*.}, leftmargin=*, itemsep=2pt]
  \item Does an SNN with online, reward-modulated STDP plasticity recover navigation
        performance significantly faster than a frozen DNN when the environment changes
        mid-episode?
  \item Does FPGA-parallel synaptic update logic reduce weight-update latency
        compared to sequential CPU/GPU computation, and at what network sizes does
        this advantage actually appear?
  \item What is the energy-efficiency tradeoff (SynOps vs.\ multiply-accumulate
        operations) between the SNN and DNN approaches at equal task performance?
\end{enumerate}
\end{keybox}

% =========================================================
\section{Background}

\subsection{Memristive Devices and Synaptic Emulation}

A memristor is a two-terminal passive element whose resistance $R(t)$ is a function of
the history of current through it \cite{strukov2008}. Its conductance $g = 1/R$ evolves
as:
\begin{equation}
  \frac{dg}{dt} = f(g, V) = \alpha\, V(t)\,\bigl[g_{\max} - g\bigr] \cdot \mathbf{1}_{V>0}
                           - \beta\, V(t)\,\bigl[g - g_{\min}\bigr] \cdot \mathbf{1}_{V<0}
  \label{eq:memristor}
\end{equation}
where $\alpha,\beta$ are potentiation/depression rates, and $[g_{\min}, g_{\max}]$
bounds the conductance. This naturally implements bounded, nonlinear, state-dependent
weight updates---the same qualitative properties (boundedness, history-dependence,
locality) that we model in software and reproduce \emph{structurally} on FPGA. The
discrete-time digital implementation does not reproduce the continuous-time
differential dynamics of Eq.~\eqref{eq:memristor}, and this distinction is maintained
throughout the thesis.

\subsection{Leaky Integrate-and-Fire Neuron}

Each neuron in the SNN follows the LIF model:
\begin{equation}
  \tau_m \frac{d\vmem(t)}{dt} = -\vmem(t) + R_m\,I_{\mathrm{syn}}(t)
  \label{eq:lif}
\end{equation}
where $\tau_m$ is the membrane time constant, $R_m$ is membrane resistance, and
$I_{\mathrm{syn}}$ is the total synaptic input current. A spike is emitted and the
membrane is reset when $\vmem(t) \geq \vth$:
\begin{equation}
  s(t) = \begin{cases} 1 & \text{if } \vmem(t) \geq \vth \\ 0 & \text{otherwise} \end{cases},
  \qquad \vmem \leftarrow V_{\mathrm{reset}} \text{ after spike.}
  \label{eq:spike}
\end{equation}
The synaptic input is filtered through an exponential kernel:
\begin{equation}
  I_{\mathrm{syn}}(t) = \sum_j w_j \sum_{t_j^f} \exp\!\left(-\frac{t - t_j^f}{\tau_s}\right)
  \label{eq:syn}
\end{equation}
where $t_j^f$ are the firing times of presynaptic neuron $j$ and $w_j$ is its synaptic
weight.

\subsection{Spike-Timing-Dependent Plasticity (STDP)}

STDP is the biologically observed learning rule in which synaptic strength is modulated
by the relative timing of pre- and post-synaptic spikes \cite{bi1998}. The weight
update for synapse $j$ is:
\begin{equation}
  \Delta w_j = \begin{cases}
    A_+\,\exp\!\left(-\dfrac{\Delta t}{\tau_+}\right) & \text{if } \Delta t > 0 \quad
      \text{(pre before post: LTP)} \\[8pt]
    -A_-\,\exp\!\left(\phantom{-}\dfrac{\Delta t}{\tau_-}\right) & \text{if } \Delta t < 0 \quad
      \text{(post before pre: LTD)}
  \end{cases}
  \label{eq:stdp}
\end{equation}
where $\Delta t = t_{\mathrm{post}} - t_{\mathrm{pre}}$, and $A_+, A_-, \tau_+, \tau_-$
are positive constants controlling potentiation and depression amplitudes and time
windows. Weights are bounded: $w_j \in [w_{\min},\, w_{\max}]$, linking directly to the
memristor conductance model in Eq.~\eqref{eq:memristor}.

\begin{notebox}
  \textbf{Link to memristor:} Eq.~\eqref{eq:stdp} is a discretized, spike-triggered
  version of Eq.~\eqref{eq:memristor}. A positive voltage pulse (pre fires before post)
  increases conductance (LTP); a negative pulse does the reverse (LTD). This is why the
  FPGA implementation of the weight-update crossbar is structurally an STDP
  accelerator whose topology mirrors a memristive crossbar---it is described as
  \emph{memristor-inspired} rather than a memristor emulator throughout this proposal.
\end{notebox}

\subsection{The Credit-Assignment Problem: Reward-Modulated STDP as the Primary Rule}
\label{sec:credit}

A central design decision of this thesis follows directly from a limitation of
Eq.~\eqref{eq:stdp}. Pure STDP is \emph{Hebbian}: it strengthens or weakens synapses
based only on the relative timing of pre- and post-synaptic spikes, with no reference
to task outcome. There is no term in the rule that encodes ``this spike pattern
preceded a collision, suppress it.'' Consequently, once an SNN is released from
surrogate-gradient pretraining into pure online STDP, weight drift is driven by the
correlation structure of the input spike trains---not by navigation success.
Surrogate-gradient pretraining establishes a good starting point, but pure STDP has no
mechanism to keep the network near that point in a \emph{task-relevant} direction. This
is precisely why Lobov et al.\ \cite{lobov2020} and, more explicitly, Juarez-Lora
et al.\ \cite{juarez2022} use \emph{reward-modulated} STDP (R-STDP), in which a global
reward/error signal gates an eligibility trace before it is consolidated into
$\Delta w_{ij}$.

Because the central hypothesis (H1) is a claim about \emph{task-relevant} recovery,
this thesis adopts \textbf{R-STDP as the primary learning rule}, and retains pure
Hebbian STDP as an ablation used to demonstrate exactly why reward modulation is
necessary:

\begin{enumerate}[leftmargin=*, itemsep=3pt]
  \item \textbf{(Primary plan) R-STDP.} An eligibility-trace buffer $e_{ij}(t)$ per
        synapse is updated by the Hebbian term of Eq.~\eqref{eq:stdp} on every spike
        pair, and a scalar reward/collision signal $r(t)$ (e.g.\ $r=-1$ on collision,
        small positive shaping reward for progress toward goal) is broadcast to all
        synapses:
        \begin{equation}
          \frac{de_{ij}}{dt} = -\frac{e_{ij}}{\tau_e} + \Delta w_{ij}^{\text{STDP}},
          \qquad
          \Delta w_{ij} = \eta\, r(t)\, e_{ij}(t).
          \label{eq:rstdp}
        \end{equation}
        The FPGA design contains the per-synapse $\Delta t$ array and LUT; R-STDP adds
        one eligibility-trace register and one scalar multiply per synapse, which is a
        modest hardware extension.
  \item \textbf{(Ablation) Pure STDP.} The same network with the reward gate removed
        ($\Delta w_{ij} = \Delta w_{ij}^{\text{STDP}}$ applied directly). This ablation
        is expected \emph{not} to recover the navigation policy in a task-relevant way;
        any residual effect is a \emph{representation-level} change (e.g.\ sensory-layer
        receptive-field drift in response to the post-shift input statistics) rather
        than policy adaptation. Contrasting R-STDP against this ablation is what
        attributes recovery specifically to reward-gated plasticity.
\end{enumerate}

\begin{riskbox}{Risk: STDP Credit Assignment}
If reward modulation were omitted, H1 would risk failing for reasons unrelated to
memristor-style plasticity (i.e.\ because Hebbian drift is task-agnostic), which would
muddy the thesis conclusion regardless of how well the FPGA or simulation work is
executed. Adopting R-STDP as the primary rule removes this confound. The
eligibility-trace and reward-broadcast machinery is small in both software and hardware
(one register + one multiply per synapse), so it is committed to from the start rather
than deferred; the pure-STDP ablation is cheap because it is the same code path with
the reward gate disabled.
\end{riskbox}

% =========================================================
\section{System Architecture}

The full system has three layers: simulation environment, neural controller, and FPGA
co-processor.

\begin{figure}[htbp]
    \centering
    % Replace with the exact filename you uploaded
    \includegraphics[width=0.6\textwidth]{mermaid-diagram.png}
    \caption{System Architecture showing the closed-loop control and FPGA co-processor.}
    \label{fig:sys_arch}
\end{figure}
\subsection{Baseline and Ablation Controllers (Comparison Group)}
\label{sec:baselines}

Four controllers are used to establish a fair and rigorous comparison:

\begin{enumerate}[leftmargin=*, itemsep=3pt]
  \item \textbf{MLP-DNN (frozen):} A multi-layer perceptron trained offline via
        supervised imitation learning on expert trajectories. After training, weights
        are fixed. This isolates the effect of online plasticity.
  \item \textbf{MLP-DNN (Online-Updating):} An MLP that performs continuous online
        learning via stochastic gradient descent (SGD) on a sliding window of recent
        observations. \emph{Loss target:} to avoid the two obvious failure modes ---
        (i) pure self-supervision on the network's own past actions, which collapses
        to a fixed point with no corrective signal, and (ii) continued access to
        expert labels during Phase~2, which would give this baseline an unfair
        advantage the SNN does not have --- the online loss is a one-step
        TD-bootstrapped value/advantage target computed from the robot's own
        collision and progress-to-goal signals (the same signals used for the
        R-STDP reward broadcast in Section~\ref{sec:credit}). This keeps the
        comparison fair: both adapting controllers learn from the same
        environment-derived feedback, differing only in \emph{how} that feedback is
        consolidated into weight changes.
  \item \textbf{SNN (frozen, ablation):} The same LIF-SNN architecture and
        surrogate-gradient-pretrained weights as the proposed controller, but with
        plasticity \emph{disabled} during evaluation. This isolates the contribution of
        spiking/LIF temporal dynamics (which alone may provide some robustness via
        low-pass filtering, per H4) from the contribution of online plasticity
        (H1). Without this ablation, any advantage of the SNN--R-STDP controller cannot
        be attributed to plasticity specifically.
  \item \textbf{RL Policy (PPO):} A proximal policy optimization agent with an MLP
        policy network, trained in PyBullet. This is the strongest fixed-policy
        baseline and represents the current state of practice for learned robotic
        control. Its policy is frozen post-training for this comparison, so it is not
        included in the recovery-time comparison but serves as an absolute
        performance ceiling for Phase~1.
\end{enumerate}

\begin{notebox}
  \textbf{Why not CNN?} CNNs are designed for grid-structured spatial data (images).
  The navigation task uses LiDAR distance readings and pose vectors---a flat,
  non-spatial observation space better suited to an MLP. Using CNN as the baseline
  would be testing the wrong inductive bias, not the right one.
\end{notebox}

\subsection{Proposed Controller: LIF-SNN with R-STDP}

\begin{enumerate}[leftmargin=*, itemsep=3pt]
  \item \textbf{Input encoding:} LiDAR readings and velocity are encoded as spike trains
        using rate coding (firing rate proportional to stimulus intensity) for distance
        channels and time-to-first-spike (TTFS) coding for urgency signals.
  \item \textbf{Network:} A fully connected LIF-SNN (3--4 layers, $\sim$512--1024
        neurons) implemented in SNNTorch \cite{eshraghian2023}. Initial weights are
        pretrained via surrogate-gradient backpropagation to match the MLP-DNN, then
        \emph{released} to online R-STDP (Section~\ref{sec:credit}) adaptation during
        evaluation.
  \item \textbf{Output decoding:} Population vote over output neurons, one population
        per discrete action (left, right, forward, brake).
  \item \textbf{Online adaptation:} At each timestep, spike times and the scalar reward
        signal are sent to the FPGA, which computes $\Delta w_{ij}$ in parallel across
        all synapses and returns updated weights before the next action.
\end{enumerate}

\subsection{FPGA Synaptic Update Coprocessor}

The key architectural claim of this thesis is that the STDP weight-update step has the
same topology and update locality as a memristor crossbar, and that this parallelism
can be hardware-exploited. On an ideal crossbar, $N \times M$ synapses update
\emph{simultaneously} in $O(1)$ time. On a CPU, this is an $O(N \cdot M)$ sequential
(or vectorized-but-still-bandwidth-bound) loop. An FPGA sits between these extremes:
it can update many synapses per cycle, but---unlike an analog crossbar---it is limited
by the number of physical DSP/multiplier resources available, so true $O(1)$
parallelism only holds up to that resource count.

The FPGA implementation (targeting a Xilinx Zynq-7000 or Intel Cyclone~V) consists of:

\begin{itemize}[leftmargin=*, itemsep=2pt]
  \item \textbf{Spike-time registers:} One register per neuron, holding the most recent
        spike timestamp $t_j^f$ in fixed-point format (16-bit, 0.1\,ms resolution).
  \item \textbf{$\Delta t$ computation array:} Parallel subtractors computing
        $\Delta t_{ij} = t_i^{\mathrm{post}} - t_j^{\mathrm{pre}}$ for a tile of
        synapses simultaneously (see resource discussion below).
  \item \textbf{Exponential lookup table (LUT):} Pre-computed STDP kernels stored in
        BRAM. The STDP window ($\tau_+,\tau_- \approx 20$--$40$\,ms) combined with
        0.1\,ms timestamp resolution spans 200--400 raw ticks; these are mapped onto
        64 LUT bins via \emph{logarithmic} (not linear) quantization of
        $|\Delta t|$, since the exponential kernel in Eq.~\eqref{eq:stdp} varies most
        steeply for small $|\Delta t|$ and is nearly flat for large $|\Delta t|$.
        This bin layout is validated against the floating-point kernel in simulation
        before synthesis (target: $<2\%$ relative error in $\Delta w$).
  \item \textbf{Weight update array:} Parallel multiply-accumulators applying the
        LTP/LTD rule and the eligibility-trace update of Eq.~\eqref{eq:rstdp}, and
        clamping to $[w_{\min}, w_{\max}]$.
  \item \textbf{AXI4-Stream (DMA) interface:} Implements Direct Memory Access for
        burst-transferring entire vectors of spike states between CPU RAM and the
        FPGA, avoiding the per-register transfer overhead of AXI4-Lite.
\end{itemize}

\begin{riskbox}{Risk: Resource-Constrained ``Parallelism''}
A fully connected $512\times512$ synapse layer requires on the order of $2.6\times10^5$
parallel multiply-accumulate units for true one-cycle updates. A Zynq-7020 provides
only 220 DSP slices. The realistic design therefore \textbf{time-multiplexes} a bank
of processing elements (PEs) across synapse tiles rather than instantiating one PE per
synapse---this is not a true $O(1)$ crossbar, and the proposal should not claim it is.
Additionally, the AXI-DMA round trip per timestep has a fixed latency floor that, for
small networks, can dominate the actual update compute. Consequently:
\begin{itemize}[leftmargin=*, itemsep=1pt]
  \item The 10--50$\times$ speedup claim (H3) applies to the \emph{inner
        weight-update kernel only}, not the end-to-end control loop including DMA
        transfer---this hedge is made explicit as a named metric in
        Section~\ref{sec:metrics}.
  \item Phase~4 includes a \textbf{network-size sweep} (e.g.\ 64, 128, 256, 512
        synapses-per-layer) reporting kernel time and end-to-end (including DMA)
        time for both CPU and FPGA, so the crossover point at which FPGA acceleration
        becomes worthwhile is an empirical result rather than an assumption. At small
        sizes, CPU/GPU vectorized BLAS may simply win, and that is a valid finding.
\end{itemize}
\end{riskbox}

This architecture is \emph{structurally analogous} to a memristive crossbar: each PE
corresponds to (a time-multiplexed share of) one synapse, spike arrival is the
analog of a voltage pulse, and the weight register is the analog of the conductance
state $g$. Unlike Eq.~\eqref{eq:memristor}, the FPGA operates in discrete time on
quantized spike events and is resource-bounded rather than truly parallel; the term
``memristor-inspired'' is used throughout to reflect this.

% =========================================================
\section{Experimental Setup}

\subsection{Simulation Environment: PyBullet}

All experiments run in PyBullet \cite{coumans2021}, a real-time physics simulation
library. The scenario:

\begin{itemize}[leftmargin=*, itemsep=2pt]
  \item \textbf{Robot:} Differential-drive ground vehicle with 360° LiDAR (32 beams,
        8\,m range) and an IMU. The observation space is deliberately kept flat and
        non-spatial (LiDAR ranges + pose/velocity) to justify the MLP inductive bias
        (Section~\ref{sec:baselines}); no camera is used.
  \item \textbf{Environment:} Enclosed arena ($10 \times 10$\,m) with randomized static
        obstacles (boxes, cylinders) and dynamic obstacles (moving at 0.3--1.0\,m/s).
  \item \textbf{Task:} Reach a goal position without collision. Episode length: 60\,s.
  \item \textbf{Environment shift:} At $t = 30$\,s, the obstacle density doubles and
        three new dynamic obstacles are injected. This is the critical test of
        adaptivity.
\end{itemize}

\subsection{The Key Experiment: Online Adaptation Under Distribution Shift}

This is the experiment that differentiates this thesis from prior work. The protocol:

\begin{enumerate}[leftmargin=*, itemsep=3pt]
  \item Run all controllers (MLP-DNN frozen, MLP-DNN online, SNN frozen
        ablation, SNN--R-STDP, and PPO) for the first 30\,s with a fixed
        environment (Phase 1).
  \item At $t = 30$\,s, inject the environment shift without stopping the episode
        (Phase 2).
  \item Measure \textbf{performance recovery time}: how many seconds until each
        controller's collision-free navigation rate returns to within 10\% of its
        Phase~1 baseline.
  \item The MLP-DNN (frozen), the SNN (frozen ablation), and the PPO policy
        \emph{cannot} adapt online. The evaluation focuses on comparing recovery
        times between the SNN--R-STDP controller and the MLP-DNN (online)
        controller, with the frozen variants providing the non-adapting reference
        level that defines ``no recovery.''
\end{enumerate}

\subsection{Evaluation Metrics}
\label{sec:metrics}

\begin{center}
\renewcommand{\arraystretch}{1.3}
\begin{tabular}{@{}p{3.2cm}p{8.2cm}p{1.6cm}@{}}
  \toprule
  \textbf{Metric} & \textbf{Definition} & \textbf{Unit} \\
  \midrule
  Response latency & $T_{\mathrm{resp}} = T_{\mathrm{action}} - T_{\mathrm{sensor}}$ & ms \\
  Energy efficiency & Synaptic operations: $\SynOps = \sum_t \sum_{ij} s_i(t)\cdot\mathbf{1}_{w_{ij}\neq 0}$ & ops/step \\
  Adaptation recovery & $\tau_{\mathrm{rec}}$: seconds to recover 90\% task performance after shift & s \\
  Robustness & Success rate under additive Gaussian noise $\sigma \in \{0, 0.1, 0.3, 0.5\}$ & \% \\
  FPGA kernel speedup & $S_{\mathrm{kernel}} = T_{\mathrm{CPU}} / T_{\mathrm{FPGA}}$ for the weight-update kernel only & $\times$ \\
  FPGA end-to-end speedup & $S_{\mathrm{e2e}} = T_{\mathrm{CPU}} / (T_{\mathrm{FPGA}} + T_{\mathrm{DMA}})$, reported per network size & $\times$ \\
  \bottomrule
\end{tabular}
\end{center}

\vspace{0.5em}
\begin{notebox}
  \textbf{Why SynOps and not $\sum|\Delta w|$?} SynOps counts how many multiply-accumulate
  operations the network actually performs per timestep, accounting for spike sparsity.
  It is the standard neuromorphic efficiency metric \cite{davies2018} and maps directly
  to hardware energy because neurons that do not spike consume no compute. $\sum|\Delta w|$
  has no physical interpretation and cannot be compared across architectures.
\end{notebox}

\subsection{Software Stack}

\begin{itemize}[leftmargin=*, itemsep=2pt]
  \item \textbf{SNN framework:} SNNTorch \cite{eshraghian2023} (Python, PyTorch backend)
  \item \textbf{RL baseline:} Stable-Baselines3 (PPO) \cite{raffin2021}
  \item \textbf{Simulation:} PyBullet \cite{coumans2021}
  \item \textbf{FPGA:} Xilinx Vivado / Intel Quartus (target: Zynq-7020 or Cyclone~V)
  \item \textbf{Hardware interface:} PyNQ (Python-to-FPGA AXI DMA bridge)
  \item \textbf{Logging:} Weights \& Biases for experiment tracking
\end{itemize}

% =========================================================
\section{Memristor-Inspired Update Model: Fidelity and Scope}

To ground the software and FPGA models in physical device behavior, the Biolek
memristor model \cite{biolek2009} is used as a \emph{conceptual reference}, not as a
target for direct emulation:
\begin{equation}
  \frac{dx}{dt} = \mu_v\,\frac{R_{\mathrm{on}}}{D^2}\,i(t)\,f(x)
  \label{eq:biolek}
\end{equation}
where $x \in [0,1]$ is the normalized state variable (analogous to normalized
conductance $g$), $\mu_v$ is the ion mobility, $D$ is the device length, and $f(x)$ is
a window function that prevents the state from leaving $[0,1]$. The STDP rule in
Eq.~\eqref{eq:stdp} is a spike-triggered, discrete-time \emph{analogy} to
Eq.~\eqref{eq:biolek}: a positive $\Delta t$ corresponds qualitatively to a
potentiating pulse, and a negative $\Delta t$ to a depressing pulse. The thesis does
not claim quantitative correspondence between $\Delta w_{ij}$ and $\Delta x$ in
Eq.~\eqref{eq:biolek}; the correspondence is structural (bounded state, history
dependence, locality of update).

The three claimed memristive \emph{properties} are reproduced structurally as follows
(software model and FPGA emulation columns describe analogous, not identical,
mechanisms):

\begin{center}
\renewcommand{\arraystretch}{1.3}
\begin{tabular}{@{}lll@{}}
  \toprule
  \textbf{Memristive Property} & \textbf{Software Model} & \textbf{FPGA Implementation} \\
  \midrule
  State retention & Weight registers persist between steps & BRAM weight array \\
  Nonlinearity \& saturation & $w \in [w_{\min}, w_{\max}]$ clamp & Saturating adder \\
  Stochasticity & Additive noise $\epsilon \sim \mathcal{N}(0, \sigma_w^2)$ & LFSR noise unit \\
  \bottomrule
\end{tabular}
\end{center}

% =========================================================
\section{Expected Results and Hypotheses}

\begin{keybox}{Hypotheses}
\begin{enumerate}[label=\textbf{H\arabic*.}, leftmargin=*, itemsep=4pt]
  \item \textbf{Adaptation:} The SNN--R-STDP controller will recover 90\% task
        performance after the environment shift within $\tau_{\mathrm{rec}} < 5$\,s,
        outperforming both the frozen MLP-DNN and the online-updating MLP-DNN. The
        pure-STDP ablation is expected \emph{not} to achieve task-relevant recovery;
        any effect it shows is confined to a measurable reorganization of sensory-layer
        representations (quantified via receptive-field drift), while the frozen-SNN
        ablation shows no such reorganization. This contrast attributes recovery
        specifically to reward-gated plasticity rather than to Hebbian drift or LIF
        dynamics alone.
  \item \textbf{Energy:} SNN SynOps will be $3$--$10\times$ lower than DNN
        multiply-accumulate operations at equal pre-shift performance, due to spike
        sparsity.
  \item \textbf{Latency:} FPGA parallel weight update via AXI-DMA will be
        $10$--$50\times$ faster than CPU computation for the weight-update
        \emph{kernel} alone, with the crossover point (network size at which FPGA
        beats CPU/GPU end-to-end) determined empirically by the Phase~4 sweep.
  \item \textbf{Robustness:} SNN success rate will degrade more gracefully under
        increasing sensor noise $\sigma$, due to the temporal integration inherent in
        LIF dynamics acting as a low-pass filter. The frozen-SNN ablation isolates
        whether this robustness comes from LIF dynamics alone or requires plasticity.
\end{enumerate}
\end{keybox}

The hypotheses are intentionally falsifiable. If H1 does not hold even for R-STDP, the
thesis conclusion is that conventional DNNs with fast online retraining are a better
practical strategy than online reward-modulated plasticity for this task class---which
is itself a valuable finding.

% =========================================================
\section{Project Timeline and Budget}
\label{sec:timeline}

\subsection{Phased Timeline (6--9 months)}

\begin{center}
\renewcommand{\arraystretch}{1.3}
\begin{tabular}{@{}clp{8.5cm}@{}}
  \toprule
  \textbf{Phase} & \textbf{Duration} & \textbf{Deliverables} \\
  \midrule
  1 & Weeks 1--4  & Literature review finalized; PyBullet environment built; MLP-DNN
                    baselines (frozen and online, with TD-based online loss) trained
                    and evaluated \\
  2 & Weeks 5--11 & SNN implemented in SNNTorch; surrogate-gradient pretraining;
                    frozen-SNN ablation evaluated; \textbf{pilot experiment}: R-STDP
                    vs.\ pure-STDP vs.\ frozen-SNN recovery under the shift (the
                    go/no-go gate before hardware) \\
  3 & Weeks 12--17 & PPO baseline trained; full comparative experiment across all
                     metrics; environment-shift protocol run for all five
                     controllers; robustness and energy sweeps (H2, H4) \\
  4 & Weeks 18--24 & FPGA R-STDP accelerator designed and synthesized (resource-aware,
                     time-multiplexed PE design); LUT bin layout validated against the
                     software golden reference; hardware-in-loop timing measurements;
                     network-size sweep for kernel and end-to-end FPGA vs.\ CPU/GPU
                     speedup \\
  5 & Weeks 25--30 & Thesis writing; figure generation; revision; presentation prep \\
  \bottomrule
\end{tabular}
\end{center}

\subsection{Budget Estimate}

\begin{center}
\renewcommand{\arraystretch}{1.3}
\begin{tabular}{@{}lrc@{}}
  \toprule
  \textbf{Item} & \textbf{Cost (EUR)} & \textbf{Notes} \\
  \midrule
  Xilinx Zynq-7020 SoC board (e.g.\ Digilent Arty Z7-20) & 180--220 &
    AXI + ARM core; PyNQ compatible \\
  \textit{Alternative:} Intel Cyclone V (DE10-Nano) & 120--150 &
    Budget option; HPS for Python side \\
  LiDAR sensor (RPLIDAR A1M8, if physical demo) & 100--130 &
    Optional; simulation covers this \\
  Differential-drive robot chassis (if physical demo) & 80--120 &
    e.g.\ Waveshare RaspiCar or similar \\
  GPU cloud compute (Vast.ai, 50 h at \textasciitilde\texteuro1.2/h) & 60--80 &
    For PPO training and SNN pretraining \\
  Miscellaneous (cables, SD cards, USB hubs) & 20--30 & \\
  \midrule
  \textbf{Total (simulation-only + FPGA)} & \textbf{260--330} & \\
  \textbf{Total (FPGA + physical robot)} & \textbf{440--580} & \\
  \midrule
  \multicolumn{3}{l}{\textit{Both scenarios are well within a 1000--1500\,EUR budget,
    leaving substantial headroom.}} \\
  \bottomrule
\end{tabular}
\end{center}

% =========================================================
\section{Related Work and Further Reading}

The following papers and resources define the landscape this thesis operates within.
They are grouped by sub-topic for targeted reading.

\subsection*{Memristive Devices and Neuromorphic Hardware}
\begin{itemize}[leftmargin=*, itemsep=2pt]
  \item Strukov et al.\ (2008), \textit{Nature}---``The missing memristor found.''
        Foundational paper establishing memristors as circuit elements \cite{strukov2008}.
  \item Biolek et al.\ (2009)---SPICE memristor model with window functions \cite{biolek2009}.
  \item Merolla et al.\ (2014), \textit{Science}---IBM TrueNorth: 4096 cores,
        1 million neurons, 256 million synapses at 70\,mW \cite{merolla2014}.
  \item Davies et al.\ (2018), \textit{IEEE Micro}---Intel Loihi: on-chip STDP,
        first programmable neuromorphic chip with online learning \cite{davies2018}.
  \item Shooshtari et al.\ (2025), \textit{Advanced Intelligent Systems}---Review of
        memristors as the hardware bridge between in-memory computing and SNNs.
\end{itemize}

\subsection*{Spiking Neural Networks: Theory and Training}
\begin{itemize}[leftmargin=*, itemsep=2pt]
  \item Maass (1997), \textit{Neural Networks}---Networks of spiking neurons as the
        third generation of neural network models \cite{maass1997}.
  \item Bi \& Poo (1998), \textit{J. Neuroscience}---Original STDP experimental
        measurement in hippocampal neurons \cite{bi1998}.
  \item Lee et al.\ (2016), \textit{Front. Neuroscience}---Training deep SNNs via
        surrogate gradients; the method used for SNN pretraining here.
  \item Eshraghian et al.\ (2023), \textit{Proc. IEEE}---SNNTorch framework paper;
        the primary software tool for this thesis \cite{eshraghian2023}.
\end{itemize}

\subsection*{Neuromorphic Robotics: Most Directly Relevant}
\begin{itemize}[leftmargin=*, itemsep=2pt]
  \item Lobov et al.\ (2020), \textit{Front. Neuroscience}---STDP-controlled mobile
        robot with memristive devices; nearest predecessor to this work \cite{lobov2020}.
  \item Juarez-Lora et al.\ (2022), \textit{Front. Neurorobotics}---R-STDP for robotic
        arm control under changing joint friction; demonstrates online adaptation and
        motivates the reward-modulated rule adopted in Section~\ref{sec:credit}
        \cite{juarez2022}.
  \item Abdelrahman et al.\ (2024/2025), \textit{Int. J. Robotics Research}---Convolutional
        SNN for obstacle avoidance on a Kinova Gen3 arm; tested in Gazebo and real world.
  \item Stroobants et al.\ (2024), \textit{arXiv:2411.13945}---First full neuromorphic
        attitude control SNN deployed on a real Crazyflie quadrotor.
  \item Supic \& Stewart (2022), \textit{ACM}---SNN adaptive control for robot arms;
        15\% faster task completion with online learning over fixed policies.
\end{itemize}

\subsection*{FPGA Neuromorphic Acceleration}
\begin{itemize}[leftmargin=*, itemsep=2pt]
  \item Qiao et al.\ (2015), \textit{Front. Neuroscience}---ROLLS neuromorphic chip
        with embedded FPGA interface for spike routing; architectural reference
        \cite{qiao2015}.
  \item Neil \& Liu (2014), \textit{IEEE BioCAS}---Minitaur: FPGA-based SNN simulator
        with on-chip STDP; close analog to the FPGA module proposed here, including
        its resource-multiplexing approach.
  \item Prajapati et al.\ (2025), \textit{arXiv:2507.20998}---Memristive SNN with
        supervised in-situ STDP; most recent circuit-level implementation reference.
\end{itemize}

\subsection*{Reinforcement Learning Baseline}
\begin{itemize}[leftmargin=*, itemsep=2pt]
  \item Schulman et al.\ (2017), \textit{arXiv:1707.06347}---PPO algorithm used for
        the RL baseline policy.
  \item Raffin et al.\ (2021), \textit{JMLR}---Stable-Baselines3 library
        \cite{raffin2021}.
\end{itemize}

% =========================================================
\section{Conclusion}

This thesis proposes a concrete, falsifiable comparative study of neuromorphic vs.\
conventional robotic control, grounded in three interlocking contributions: (1) a
software SNN with biologically grounded, reward-modulated STDP, evaluated with
frozen-SNN and pure-STDP ablations that separate the contributions of spiking/LIF
dynamics, Hebbian drift, and reward-gated online plasticity under environment shift,
(2) an FPGA coprocessor whose update topology is structurally analogous to a memristive
crossbar---explicitly scoped as resource-constrained and time-multiplexed rather than
truly $O(1)$---and (3) a quantitative evaluation framework using standard metrics
(SynOps, recovery time, kernel and end-to-end FPGA speedup) that allows direct
comparison with the existing literature.

The project is achievable within a 6--9 month timeline and a budget of 440--580\,EUR
including physical hardware, or under 350\,EUR for a simulation-only study. The FPGA
component is not decorative---it is the hardware argument for why memristor-inspired,
in-memory synaptic computation is worth building, even within the resource limits of a
commodity FPGA: it makes the parallel-update structure of synaptic learning
physically instantiable without fabricating nanoscale devices.

% =========================================================
\begin{thebibliography}{99}
\small

\bibitem{strukov2008}
D.~B. Strukov, G.~S. Snider, D.~R. Stewart, and R.~S. Williams,
``The missing memristor found,''
\textit{Nature}, vol.~453, pp.~80--83, 2008.

\bibitem{bi1998}
G.~Q. Bi and M.~M. Poo,
``Synaptic modifications in cultured hippocampal neurons: dependence on spike timing,
synaptic strength, and postsynaptic cell type,''
\textit{Journal of Neuroscience}, vol.~18, no.~24, pp.~10464--10472, 1998.

\bibitem{maass1997}
W.~Maass,
``Networks of spiking neurons: The third generation of neural network models,''
\textit{Neural Networks}, vol.~10, no.~9, pp.~1659--1671, 1997.

\bibitem{merolla2014}
P.~A. Merolla \textit{et al.},
``A million spiking-neuron integrated circuit with a scalable communication network
and interface,''
\textit{Science}, vol.~345, no.~6197, pp.~668--673, 2014.

\bibitem{davies2018}
M.~Davies \textit{et al.},
``Loihi: A neuromorphic manycore processor with on-chip learning,''
\textit{IEEE Micro}, vol.~38, no.~1, pp.~82--99, 2018.

\bibitem{biolek2009}
Z.~Biolek, D.~Biolek, and V.~Biolkova,
``SPICE model of memristor with nonlinear dopant drift,''
\textit{Radioengineering}, vol.~18, no.~2, pp.~210--214, 2009.

\bibitem{eshraghian2023}
J.~K. Eshraghian \textit{et al.},
``SNNTorch: Accelerating neuromorphic engineering with a deep learning algorithm
library,''
\textit{Proceedings of the IEEE}, vol.~111, no.~9, pp.~1--19, 2023.

\bibitem{lobov2020}
S.~A. Lobov \textit{et al.},
``Spatial properties of STDP in a self-learning spiking neural network enable
controlling a mobile robot,''
\textit{Frontiers in Neuroscience}, vol.~14, p.~88, 2020.

\bibitem{juarez2022}
A.~Juarez-Lora \textit{et al.},
``R-STDP spiking neural network architecture for motion control on a changing friction
joint robotic arm,''
\textit{Frontiers in Neurorobotics}, vol.~16, 2022.

\bibitem{qiao2015}
N.~Qiao \textit{et al.},
``A reconfigurable on-line learning spiking neuromorphic processor comprising 256
neurons and 128K synapses,''
\textit{Frontiers in Neuroscience}, vol.~9, p.~141, 2015.

\bibitem{raffin2021}
A.~Raffin \textit{et al.},
``Stable-Baselines3: Reliable reinforcement learning implementations,''
\textit{Journal of Machine Learning Research}, vol.~22, no.~268, pp.~1--8, 2021.

\bibitem{coumans2021}
E.~Coumans and Y.~Bai,
``PyBullet, a Python module for physics simulation for games, robotics and machine
learning,'' \url{http://pybullet.org}, 2016--2021.

\end{thebibliography}

\end{document}
